Humanity is driven to understand the nature of the universe, leading to great efforts and achievements in theoretical and experimental physics.
One of these great achievements, a combined effort of many across decades and through multiple experiments, is the formulation of a framework to describe the fundamental particles and their interactions.
This framework is known as the Standard Model (SM) of particle physics.
Complete in its ability to describe and predict the observed phenomena of the particles it contains, it remains the starting point for virtually every investigation in particle physics.
Some of the most recent confirmations of the SM are the top quark observation at Fermilab in 1995 by the D0 and CDF experiments \cite{PhysRevLett.74.2626,PhysRevLett.74.2632} and the Higgs boson discovery at CERN in 2012 by the ATLAS and CMS experiments \cite{HIGG-2012-27,CMS-HIG-12-028}.
Recent measurements of the W boson mass, which initially showed a significant deviation from the SM prediction, have also been updated with new measurements that are consistent with the SM \cite{science.abk1781,atlas_2024}, further confirming that the SM is able to withstand intense scrutiny.


Despite its successes, however, the Standard Model is known to be incomplete in regards to several observations.
We know that over 85\% of the matter in the observable universe is in the form of dark matter (DM), which does not have any observed interactions with the SM and thus must be explained with new physics.
Additionally, at certain energy scales, the SM requires neutrinos to be massless, which is contradicted by neutrino oscillation observations.
The extensions to the SM and discussion needed to explain neutrino masses are beyond the scope of this thesis.
The matter-antimatter asymmetry observed in the universe is also not fully explained by the SM, requiring some new mechanism to generate the observed abundance of matter.
A strongly-coupled theory of quantum gravity is unable to be incorporated into the SM, and there is no explanation within the SM for dark energy at the scales observed in cosmology.
These shortcomings and others motivate searches for physics beyond the Standard Model (BSM).

This thesis focuses on two specific BSM searches with the ATLAS detector at the Large Hadron Collider (LHC) at CERN.
Motivated by two different approaches to explaining dark matter, the first search is for dark mesons decaying into top and bottom quarks \cite{kribs_2019}, and the second search is for dark photons from Higgs decays in the gluon-gluon fusion production mode \cite{Biswas_2022}.
Both models are explained in more detail in Chapter \ref{ch:theory}, along with an overview of the Standard Model and the motivation for BSM physics.
The production cross sections for the these signals are $\sim 10^{13}$ times smaller than the total inelastic proton-proton cross section at the LHC, so these searches require sophisticated analysis techniques to separate the signal from the overwhelming background.
Chapter \ref{ch:experimental_setup} gives an overview of the experimental methods used in these searches, including a description of the ATLAS detector and how the LHC provides an unprecedented center of mass energy for proton-proton collisions.
The LHC will enter a new phase of operation in the coming years, with the High-Luminosity LHC (HL-LHC) upgrade providing a significant increase in luminosity and the ATLAS Phase II upgrade improving the detector capabilities to handle the increased pile-up and radiation damage.
Chapter \ref{ch:hl-lhc-atlas-upgrades} gives an overview of the HL-LHC and ATLAS Phase II upgrades, and the impact they will have on the physics reach of the ATLAS experiment.
The two searches for dark mesons and dark photons are described in Chapters \ref{ch:dark_mesons} and \ref{ch:dark_photons}, respectively.
My primary contributions to the dark meson search include a truth study of the multi-lepton channel, verifying preselection yields for the all-hadronic channel, verifying the 4-d ABCD method for the all-hadronic channel, and assisting in the development of the analysis and plotting frameworks for the one-lepton channel.
As a main contributor to the dark photon search, I was heavily involved in the development and implementation of the fitting regions for the statistical analysis, validation of the background estimation methods and vertex BDT, optimization of the signal region, studies to improve the primary vertex selection, and studies of a new neural network method to reconstruct \ETmiss.
Concluding remarks and future outlook are given in Chapter \ref{ch:conclusion}.


% This chapter will discuss the Standard Model briefly, with emphasis on electroweak symmetry breaking and the Higgs mechanism, before discussing the motivation for BSM physics and two possible dark sectors in more detail.